\documentclass[14pt]{extarticle}
\usepackage{fontspec}
\usepackage{polyglossia}
\usepackage{amsmath}
\usepackage{amssymb}
\usepackage{xcolor}
\usepackage{fancyhdr}
\usepackage{graphicx}
\usepackage{listings}
\usepackage[most]{tcolorbox}
% \usepackage{setspace}
\usepackage{pifont}
\usepackage{enumitem}
\usepackage{titlesec}
\usepackage[bottom]{footmisc}
\usepackage{titling}
\usepackage{minted}
\usepackage{etoolbox}
\usepackage{array}
\usepackage{longtable} % page-breakable tables (must load before polyglossia/bidi RTL setup)
\usepackage{extsizes}
\usepackage{float}
\usepackage{booktabs}
\usepackage{hyperref}
\usepackage[onehalfspacing]{setspace}
\newif\ifwhiteblack
\ifdefined\setwhiteblack
  \setwhiteblack
\fi

\ifwhiteblack
  \usepackage[a4paper, landscape, top=1.25cm, bottom=1.25cm, left=1.25cm, right=1.25cm, headheight=1.5cm, headsep=0.5cm]{geometry}
\else
  \usepackage[a4paper, top=2.7cm, bottom=2.7cm, left=2.5cm, right=2.5cm, headheight=1.5cm, headsep=0.5cm]{geometry}
\fi
\usepackage{tikz}
\usepackage{datetime2}

\usetikzlibrary{automata, positioning, arrows.meta, arrows}
\providecommand{\LTcaptype}{table}

\newfontfamily\emoji{Segoe UI Emoji}

\pagestyle{fancy}

\setmainlanguage[numerals=western]{arabic}
\setotherlanguage{english}
\newfontfamily\arabicfont[Script=Arabic]{Calibri}
\newfontfamily\arabicfonttt[Script=Arabic]{Courier New}



\ifwhiteblack
\lstset{
  language=[Sharp]C,
  numbers=left,
  stepnumber=1,
  numbersep=8pt,
  frame=single,
  basicstyle=\ttfamily\small,
  keywordstyle=\bfseries,
  stringstyle=\ttfamily,
  commentstyle=\itshape
}
\else
\lstset{
  language=[Sharp]C,
  numbers=left,
  stepnumber=1,
  numbersep=8pt,
  frame=single,
  basicstyle=\ttfamily\small,
  keywordstyle=\color{blue},
  stringstyle=\color{red},
  commentstyle=\color{green!50!black}
}
\fi

\newif\ifdetailed
\ifdefined\setdetailed
  \setdetailed
\fi

\newif\ifwithsols
\ifdefined\setwithsols
  \setwithsols
\fi

\ifwhiteblack
  \AtBeginDocument{
    \LARGE
    \usemintedstyle{bw}
  }
\fi

% Default minted look: compact code that wraps long lines.
% Not applied in black-and-white landscape mode, which already uses \LARGE.
\ifwhiteblack\else
  \setminted{fontsize=\small, breaklines=true}
\fi

% unified tcolorboxes for articles
\ifwhiteblack
\tcbset{colback=white, colframe=black, fonttitle=\bfseries, boxrule=1pt}
\newtcolorbox{boxDef}[1][]{title={{\emoji📘} تعريف\ifx\\#1\\\else ~#1\fi :}}
\newtcolorbox{boxExercise}[1][]{title={{\emoji🧩} تمرين\ifx\\#1\\\else ~#1\fi :}}
\newtcolorbox{boxExample}[1][]{title={{\emoji📝} مثال\ifx\\#1\\\else ~#1\fi :}}
\newtcolorbox{boxNote}[1][]{title={{\emoji✨} ملاحظة\ifx\\#1\\\else ~#1\fi :}}
\newtcolorbox{boxAttention}[1][]{title={{\emoji🔔} تنبيه\ifx\\#1\\\else ~#1\fi :}}
\newtcolorbox{boxWarning}[1][]{title={{\emoji⚡} ملاحظة هامة\ifx\\#1\\\else ~#1\fi :}}
\newtcolorbox{boxSolution}[1][]{title={{\emoji✅} حل\ifx\\#1\\\else ~#1\fi :}}
\newtcolorbox{boxSymbol}[1][]{title={{\emoji🔣} رمز\ifx\\#1\\\else ~#1\fi :}}
\newtcolorbox{boxHint}[1][]{title={{\emoji💡} تلميح\ifx\\#1\\\else ~#1\fi :}}
\else
\tcbset{colback=white, colframe=black, fonttitle=\bfseries, boxrule=0.8pt}
\newtcolorbox{boxDef}[1][]{colback=blue!5!white,colframe=blue!75!black,
  title={{\emoji📘} تعريف\ifx\\#1\\\else ~#1\fi :}}
\newtcolorbox{boxExercise}[1][]{colback=cyan!5!white,colframe=cyan!70!black,
  title={{\emoji🧩} تمرين\ifx\\#1\\\else ~#1\fi :}}
\newtcolorbox{boxExample}[1][]{colback=yellow!5!white,colframe=orange!90!black,
  title={{\emoji📝} مثال\ifx\\#1\\\else ~#1\fi :}}
\newtcolorbox{boxNote}[1][]{colback=gray!10!white,colframe=black,
  title={{\emoji✨} ملاحظة\ifx\\#1\\\else ~#1\fi :}}
\newtcolorbox{boxAttention}[1][]{colback=magenta!10!white,colframe=magenta!80!black,
  title={{\emoji🔔} تنبيه\ifx\\#1\\\else ~#1\fi :}}
\newtcolorbox{boxWarning}[1][]{colback=red!5!white,colframe=red!75!black,
  title={{\emoji⚡} ملاحظة هامة\ifx\\#1\\\else ~#1\fi :}}
\newtcolorbox{boxSolution}[1][]{colback=green!5!white,colframe=green!60!black,
  title={{\emoji✅} حل\ifx\\#1\\\else ~#1\fi :}}
\newtcolorbox{boxSymbol}[1][]{colback=purple!5!white,colframe=purple!70!black,
  title={{\emoji🔣} رمز\ifx\\#1\\\else ~#1\fi :}}
\newtcolorbox{boxHint}[1][]{colback=teal!5!white,colframe=teal!60!black,
  title={{\emoji💡} تلميح\ifx\\#1\\\else ~#1\fi :}}
\fi


\ifwhiteblack
\tcbset{simplecode/.style={ colback=white, colframe=black, boxrule=1pt, arc=2pt, left=4pt,right=4pt,top=4pt,bottom=4pt}}
\else
\tcbset{simplecode/.style={ colback=gray!5, colframe=black!50, boxrule=0.4pt, arc=2pt, left=4pt,right=4pt,top=4pt,bottom=4pt}}
\fi
% Compact code boxes: tighter line spacing (the preamble uses onehalfspacing) and small vertical gaps.
% Must come AFTER the simplecode definition above.
\tcbset{simplecode/.append style={fontupper=\setstretch{0.95}, top=4pt, bottom=4pt, before skip=6pt, after skip=6pt}}
\newenvironment{boxCode}{\begin{tcolorbox}[simplecode]}{\end{tcolorbox}}

% Code cell used inside tables (keeps code left-to-right inside RTL tables).
% Use inside a fixed-width p{...} column: \begin{pycell}\begin{minted}...\end{minted}\end{pycell}
\newtcolorbox{pycell}{frame empty,colback=white,boxsep=0pt,left=2pt,right=2pt,top=2pt,bottom=2pt,width=\linewidth}

\newcolumntype{C}[1]{>{\centering\arraybackslash}p{#1}}

\makeatletter
\DTMnewdatestyle{dotted}{%
  \renewcommand{\DTMdisplaydate}[4]{\two@digits{##3}.\two@digits{##2}.##1}%
  \renewcommand{\DTMDisplaydate}{\DTMdisplaydate}%
}
\makeatother
\DTMsetdatestyle{dotted}

% redefine spaces after titles
\makeatletter
\renewcommand{\@maketitle}{%
  \begin{center}
    {\huge \bfseries \@title \par}%
    \vskip 0.2em % space between title and author
    {\large \@author \par}%
    % \vskip 0.2em % space between author and date
    % {\normalsize \@date \par}%
    {\small \hfill \begin{english}\DTMtoday\end{english} \par}%
  \end{center}
}
\makeatother

\fancyhf{} % clear default
\fancypagestyle{plain}{
  \fancyhf{}
  \fancyhead[L]{مدرسة التسامح الشاملة}
  % \fancyhead[L]{\includegraphics[height=1cm]{../../../images/logoTasamoh.png}}
  \fancyhead[R]{الأستاذ محمود اغبارية}
  \fancyfoot[C]{\thepage}
}

\fancyhead[L]{مدرسة التسامح الشاملة}
\fancyhead[R]{الأستاذ محمود اغبارية}
\fancyfoot[C]{\thepage}

\setcounter{tocdepth}{3} % only section subsection and subsubsection in TOC

\newcommand{\importpdfpage}[2]{%
    \thispagestyle{empty}
    \newgeometry{left=0cm,right=0cm,top=0cm,bottom=0cm}
    \includegraphics[page=#2,width=0.95\paperwidth,height=0.95\paperheight,keepaspectratio]{#1}
    \restoregeometry
}

\newcommand{\insertFullImg}[1]{%
    \noindent
    \makebox[\textwidth][c]{
        \includegraphics[width=0.9\paperwidth,keepaspectratio]{#1}%
    }%
}%


\ifwithsols
\title{حل ورقة عمل 3: الوراثة وتعدد الأشكال}
\else
\title{ورقة عمل 3: الوراثة وتعدد الأشكال}
\fi

\begin{document}

\maketitle
\thispagestyle{fancy}

\begin{boxDef}[مقدمة وسياق المسألة]
في شركة \textbf{"السعادة للجميع"}، يوجد 4 أنواع من الموظفين، بحيث ينتمي كل موظف إلى نوع واحد فقط:
\begin{itemize}
    \item \textbf{موظف عادي:} موظف عادي يعمل في الشركة.
    \item \textbf{عضو لجنة العمال:} موظف يشغل عضوية في نقابة/لجنة عمال الشركة.
    \item \textbf{تقني حواسيب:} موظف مسؤول عن صيانة عدد من الحواسيب (كل تقني مسؤول عن عدد محدد من الحواسيب).
    \item \textbf{مسؤول مشاريع:} موظف يدير مشاريع خاصة في الشركة، وهو مسؤول عن فريق مكوّن من 15 موظفاً يضم موظفين عاديين، وأعضاء لجنة، وتقنيين.
\end{itemize}
\end{boxDef}

\begin{boxNote}[شروط التوظيف والترقية وتحديد الرواتب]
\begin{itemize}
    \item \textbf{التوظيف والترقية:} عند التعيين الجديد يتم قبول موظفين عاديين، كما يمكن توظيف تقنيين مباشرة. وبموجب قرار مجلس إدارة الشركة، فإن أعضاء لجنة العمال ومسؤولي المشاريع تتم ترقيتهم حصراً من كادر الموظفين العاديين أو التقنيين.
    \item \textbf{سلّم الرواتب الشهري:} يتحدد راتب الموظف وفق الراتب الأساسي، والأقدمية، وعلاوات إضافية:
    \begin{itemize}
        \item الراتب الأساسي الابتدائي: $5000$ ش.ج.
        \item عن كل سنة أقدمية: علاوة قدرها $500$ ش.ج.
        \item عن كل 10 حواسيب (للتقني): علاوة قدرها $250$ ش.ج \\ (أي: $\lfloor \frac{\text{computers}}{10} \rfloor \cdot 250$).
        \item منحة إدارة خاصة (لمسؤول المشاريع): علاوة قدرها $100$ ش.ج عن كل تقني، وعلاوة قدرها $50$ ش.ج عن كل عضو لجنة عمال ضمن فريق مشروعه المكون من 15 موظفاً.
    \end{itemize}
\end{itemize}
\end{boxNote}

\section*{نظام التصويت ونقاط التأثير}
تُتخذ القرارات في الشركة عبر التصويت بمشاركة جميع الموظفين. يتم التصويت بواسطة \textbf{نقاط التأثير}، حيث يمتلك كل موظف عدداً مختلفاً من النقاط وفق نوعه وأقدميته كما هو موضح في الجدول التالي:

\begin{center}
\renewcommand{\arraystretch}{1.6}
{\small
\begin{tabular}{|p{0.22\textwidth}|p{0.72\textwidth}|}
\hline
\multicolumn{1}{|c|}{\textbf{نوع الموظف}} & \multicolumn{1}{c|}{\textbf{عدد نقاط التصويت}} \\
\hline
\textbf{موظف عادي} \newline (\textenglish{Employee}) &
يحصل على 4 نقاط لكونه موظفاً في الشركة، ونقطة إضافية عن كل سنة أقدمية في الشركة. \newline
\textbf{مثال:} موظف عادي لديه 5 سنوات أقدمية يحصل على $4 + 5 = 9$ نقاط. \\
\hline
\textbf{عضو لجنة عمال} \newline (\textenglish{UnionMember}) &
يحصل على \textbf{ضعف} عدد النقاط التي يستحقها الموظف العادي، بالإضافة إلى نقطتين عن كل سنة يشغل فيها عضوية لجنة العمال. \newline
\textbf{مثال:} عضو لجنة لديه 5 سنوات أقدمية في الشركة ومن بينها 3 سنوات في اللجنة، يحصل على: $2 \cdot 9 + 2 \cdot 3 = 18 + 6 = 24$ نقطة. \\
\hline
\textbf{تقني حواسيب} \newline (\textenglish{Technician}) &
يحصل على نقاط مثل الموظف العادي، بالإضافة إلى نقطة واحدة عن كل حاسوب تحت مسؤوليته. \newline
\textbf{مثال:} تقني لديه 5 سنوات أقدمية ومسؤول عن 4 حواسيب، يحصل على: $9 + 4 = 13$ نقطة. \\
\hline
\textbf{مسؤول مشاريع} \newline (\textenglish{Supervisor}) &
يحصل على نقاط مثل الموظف العادي، بالإضافة إلى مجموع نقاط جميع الموظفين الـ 15 الذين يشرف عليهم. \newline
\textbf{مثال:} مسؤول مشاريع لديه 5 سنوات أقدمية ومجموع نقاط الـ 15 موظفاً في فريقه هو 80 نقطة، يحصل على: $9 + 80 = 89$ نقطة. \\
\hline
\end{tabular}
}
\end{center}

ترغب إدارة الشركة في بناء نظام محوسب يحسب ويعرض تكلفة الراتب الشهري وعدد نقاط التصويت لكل موظف. لهذا الغرض، تقرر تعريف الفئات التالية:
\textenglish{Employee}، \textenglish{UnionMember}، \textenglish{Technician}، \textenglish{Supervisor}، وفئة الشركة \textenglish{Company}.

\clearpage

\section*{توصيف واجهات الفئات}

فيما يلي التوصيف الكامل للصفات والعمليات لكل فئة من الفئات المطلوب تطبيقها في النظام:

\subsection*{1. الفئة الأساسية: \textenglish{Employee}}
أمامك واجهة الفئة الأساسية \textenglish{Employee} التي تمثّل موظفاً عاماً في الشركة:

\begin{center}
\renewcommand{\arraystretch}{1.5}
{\small
\begin{tabular}{|p{0.43\textwidth}|p{0.49\textwidth}|}
\hline
\multicolumn{2}{|c|}{\textbf{\textenglish{Employee}}} \\
\hline
صفة محمية تُمثّل رقم هوية الموظف. & \textenglish{protected string id} \\
\hline
صفة محمية تُمثّل اسم الموظف. & \textenglish{protected string name} \\
\hline
صفة محمية تُمثّل أقدمية الموظف بالسنوات. & \textenglish{protected int seniority} \\
\hline
عمليّة بنائيّة تتلقّى قيمًا لصفات الموظف وتهيئها. & \textenglish{public Employee(string id, string name, int seniority)} \\
\hline
عمليّة تُعيد رقم هوية الموظف. & \textenglish{public string GetId()} \\
\hline
عمليّة تُعدّل رقم هوية الموظف. & \textenglish{public void SetId(string id)} \\
\hline
عمليّة تُعيد اسم الموظف. & \textenglish{public string GetName()} \\
\hline
عمليّة تُعدّل اسم الموظف. & \textenglish{public void SetName(string name)} \\
\hline
عمليّة تُعيد عدد سنوات أقدمية الموظف. & \textenglish{public int GetSeniority()} \\
\hline
عمليّة تُعدّل عدد سنوات أقدمية الموظف. & \textenglish{public void SetSeniority(int seniority)} \\
\hline
عمليّة تحسب وتُعيد نقاط التصويت ($4 + \text{seniority}$). & \textenglish{public virtual int GetScore()} \\
\hline
عمليّة تحسب وتُعيد الراتب الشهري ($5000 + 500 \cdot \text{seniority}$). & \textenglish{public virtual double GetSalary()} \\
\hline
عمليّة تفحص استحقاق المكافأة السنوية ($\text{seniority} \ge 5$). & \textenglish{public virtual bool YearlyBonus()} \\
\hline
عمليّة تُعيد وصفاً نصياً للموظف وبياناته. & \textenglish{public override string ToString()} \\
\hline
\end{tabular}
}
\end{center}

\clearpage

\subsection*{2. الفئة الفرعية: \textenglish{Technician}}
أمامك واجهة الفئة الفرعية \textenglish{Technician} التي ترث من \textenglish{Employee} وتُمثّل تقني الحواسيب:

\begin{center}
\renewcommand{\arraystretch}{1.5}
{\small
\begin{tabular}{|p{0.43\textwidth}|p{0.49\textwidth}|}
\hline
\multicolumn{2}{|c|}{\textbf{\textenglish{Technician : Employee}}} \\
\hline
صفة خاصّة تُمثّل عدد الحواسيب تحت مسؤوليته. & \textenglish{private int computers} \\
\hline
عمليّة بنائيّة تتلقّى قيم صفات الموظف وعدد الحواسيب، وتهيئها بالاستعانة بالفئة الأساسية (\textenglish{base}). & \textenglish{public Technician(string id, string name, int seniority, int computers)} \\
\hline
عمليّة تُعيد عدد الحواسيب. & \textenglish{public int GetComputers()} \\
\hline
عمليّة تُعدّل عدد الحواسيب. & \textenglish{public void SetComputers(int computers)} \\
\hline
عمليّة تحسب وتُعيد نقاط التصويت (نقاط الموظف العادي + نقطة لكل حاسوب). & \textenglish{public override int GetScore()} \\
\hline
عمليّة تحسب وتُعيد الراتب الشهري (راتب الموظف العادي + 250 ش.ج عن كل 10 حواسيب). & \textenglish{public override double GetSalary()} \\
\hline
عمليّة تفحص استحقاق المكافأة السنوية (أقدمية 5 سنوات و100 حاسوب على الأقل). & \textenglish{public override bool YearlyBonus()} \\
\hline
عمليّة تُعيد وصفاً نصياً للتقني وبياناته وعدد حواسيبه. & \textenglish{public override string ToString()} \\
\hline
\end{tabular}
}
\end{center}

\clearpage

\subsection*{3. الفئة الفرعية: \textenglish{UnionMember}}
أمامك واجهة الفئة الفرعية \textenglish{UnionMember} التي ترث من \textenglish{Employee} وتُمثّل عضو لجنة العمال:

\begin{center}
\renewcommand{\arraystretch}{1.5}
{\small
\begin{tabular}{|p{0.43\textwidth}|p{0.49\textwidth}|}
\hline
\multicolumn{2}{|c|}{\textbf{\textenglish{UnionMember : Employee}}} \\
\hline
صفة خاصّة تُمثّل عدد سنوات عضويته في لجنة العمال. & \textenglish{private int unionYears} \\
\hline
عمليّة بنائيّة تتلقّى قيم صفات الموظف وسنوات عضويته باللجنة، وتهيئها بالاستعانة بالفئة الأساسية (\textenglish{base}). & \textenglish{public UnionMember(string id, string name, int seniority, int unionYears)} \\
\hline
عمليّة تُعيد عدد سنوات عضويته في اللجنة. & \textenglish{public int GetUnionYears()} \\
\hline
عمليّة تُعدّل عدد سنوات عضويته في اللجنة. & \textenglish{public void SetUnionYears(int unionYears)} \\
\hline
عمليّة تحسب وتُعيد نقاط التصويت (ضعف نقاط الموظف العادي + نقطتان لكل سنة باللجنة). & \textenglish{public override int GetScore()} \\
\hline
عمليّة تُعيد الراتب الشهري (يماثل راتب الموظف العادي). & \textenglish{public override double GetSalary()} \\
\hline
عمليّة تفحص استحقاق المكافأة السنوية (تُعيد \textenglish{false} دائماً). & \textenglish{public override bool YearlyBonus()} \\
\hline
عمليّة تُعيد وصفاً نصياً لعضو اللجنة وسنوات عضويته. & \textenglish{public override string ToString()} \\
\hline
\end{tabular}
}
\end{center}

\clearpage

\subsection*{4. الفئة الفرعية: \textenglish{Supervisor}}
أمامك واجهة الفئة الفرعية \textenglish{Supervisor} التي ترث من \textenglish{Employee} وتُمثّل مسؤول المشاريع:

\begin{center}
\renewcommand{\arraystretch}{1.5}
{\small
\begin{tabular}{|p{0.43\textwidth}|p{0.49\textwidth}|}
\hline
\multicolumn{2}{|c|}{\textbf{\textenglish{Supervisor : Employee}}} \\
\hline
صفة خاصّة تُمثّل مصفوفة تضم 15 موظفاً في فريق مشروعه. & \textenglish{private Employee[] team} \\
\hline
عمليّة بنائيّة تتلقّى قيم صفات الموظف ومصفوفة فريقه، وتهيئها بالاستعانة بالفئة الأساسية (\textenglish{base}). & \textenglish{public Supervisor(string id, string name, int seniority, Employee[] team)} \\
\hline
عمليّة تُعيد مصفوفة موظفي الفريق. & \textenglish{public Employee[] GetTeam()} \\
\hline
عمليّة تُعدّل مصفوفة موظفي الفريق. & \textenglish{public void SetTeam(Employee[] team)} \\
\hline
عمليّة تحسب وتُعيد نقاط التصويت (نقاط الموظف العادي + مجموع نقاط كافة أعضاء فريقه). & \textenglish{public override int GetScore()} \\
\hline
عمليّة تحسب وتُعيد الراتب الشهري (راتب الموظف العادي + 100 ش.ج لكل تقني و50 ش.ج لكل عضو لجنة بفريقه). & \textenglish{public override double GetSalary()} \\
\hline
عمليّة تفحص استحقاق المكافأة السنوية (تُعيد \textenglish{false} دائماً). & \textenglish{public override bool YearlyBonus()} \\
\hline
عمليّة تُعيد وصفاً نصياً لمسؤول المشاريع وعدد أعضاء فريقه. & \textenglish{public override string ToString()} \\
\hline
\end{tabular}
}
\end{center}

\clearpage

\subsection*{5. فئة إدارة الشركة: \textenglish{Company}}
أمامك واجهة فئة الشركة \textenglish{Company} التي تدير موظفي الشركة وتحتوي على قائمة الموظفين:

\begin{center}
\renewcommand{\arraystretch}{1.5}
{\small
\begin{tabular}{|p{0.43\textwidth}|p{0.49\textwidth}|}
\hline
\multicolumn{2}{|c|}{\textbf{\textenglish{Company}}} \\
\hline
ثابت عام يُمثّل الحد الأقصى للموظفين في الشركة (100). & \textenglish{public const int MAX = 100} \\
\hline
صفة خاصّة تُمثّل مصفوفة موظفي الشركة. & \textenglish{private Employee[] employees} \\
\hline
صفة خاصّة تُمثّل عدد الموظفين الحاليين الفعليين في الشركة. & \textenglish{private int current} \\
\hline
عمليّة بنائيّة تُنشئ شركة جديدة خالية من الموظفين (تُهيئ المصفوفة بحجم \textenglish{MAX} والعدد الحالي 0). & \textenglish{public Company()} \\
\hline
عمليّة تتلقّى موظفاً وتضيفه لمصفوفة موظفي الشركة إذا كان هناك متسع وتُعيد \textenglish{true}، وإلا تُعيد \textenglish{false}. & \textenglish{public bool AddEmployee(Employee e)} \\
\hline
عمليّة تحسب وتُعيد مجموع الرواتب الشهرية لجميع موظفي الشركة. & \textenglish{public double TotalSalary()} \\
\hline
\end{tabular}
}
\end{center}

\clearpage

\section*{المهام والمطلوب}

حل البنود التالية بأفضل صورة تراعي مبادئ البرمجة كائنية التوجه:

\begin{enumerate}[label=\textbf{\alph*.} , itemsep=2.5em]

    \item \textbf{تطبيق فئات الموظفين والعمليات البنائية:} \\
    اكتب تطبيقاً كاملاً لتعريف وخصائص والعمليات البنائية لفئات الموظفين الأربع: \newline
    \textenglish{Employee}، \textenglish{Technician}، \textenglish{UnionMember}، و \textenglish{Supervisor}، \newline
    مع الاستفادة من مبدأ الوراثة واستدعاء العملية البنائية للفئة الأساسية (\textenglish{base}) كما هو موضح في الجداول أعلاه.

    \item \textbf{العملية \textenglish{ToString}:} \\
    طبّق العملية \textenglish{ToString} في كل فئة من فئات الموظفين الأربع، بحيث تعيد وصفاً نصياً شاملاً للموظف وقيمه، مع تجنب تكرار الكود عبر استدعاء \textenglish{base.ToString()}.

    \item \textbf{العمليتان \textenglish{GetScore} و \textenglish{GetSalary}:} \\
    طبّق العمليتين \textenglish{GetScore} و \textenglish{GetSalary} في كل فئة من فئات الموظفين الأربع، مع الاستفادة من تعدد الأشكال واستدعاء عمليات الفئة الأساسية لتجنب تكرار العمليات الحسابية.

    \item \textbf{عملية المكافأة السنوية (\textenglish{YearlyBonus}):} \\
    حققت الشركة أرباحاً ممتازة وترغب في توزيع مكافأة سنوية للموظفين. حددت الشركة استحقاق المكافأة للموظفين العاديين والتقنيين \textbf{فقط} وفق الشروط التالية:
    \begin{itemize}
        \item أقدمية 5 سنوات على الأقل في الشركة.
        \item بالنسبة للتقنيين: مسؤولية صيانة 100 حاسوب على الأقل بالإضافة لشرط الأقدمية.
    \end{itemize}
    طبّق العملية \textenglish{YearlyBonus} في الفئات المناسبة بحيث تفحص هل يستحق الموظف المكافأة السنوية، وبما يتوافق مع مبادئ الوراثة وتعدد الأشكال وهندسة البرمجيات.
\clearpage
    \item \textbf{تطبيق فئة الشركة (\textenglish{Company}):} \\
    اكتب تطبيقاً كاملاً لفئة الشركة \textenglish{Company} يشمل الثابت، الخصائص، العملية البنائية، وعمليتي \textenglish{AddEmployee} و \textenglish{TotalSalary} الموضحة في جدول الفئة أعلاه. وضّح علاقة هذه الفئة بهرمية الفئات.

    \item \textbf{عملية فحص قبول الاقتراح (\textenglish{IsAccepted}):} \\
    طُرح اقتراح رسمي في الشركة لتغيير اسمها من "السعادة للجميع" إلى \textbf{"السعادة الآن وفي المستقبل"}. \\
    صوّت جميع الموظفين العاديين والتقنيين \textbf{تأييداً} للاقتراح (مع)، بينما عارض جميع أعضاء لجنة العمال ومسؤولي المشاريع هذا الاقتراح (\textbf{ضد}). \\
    يُعتمد القرار فقط إذا كان مجموع نقاط المؤيدين \textbf{أكبر تماماً} من مجموع نقاط المعارضين. \\
    أمامك تعريف عملية عامة وثابتة:
\begin{boxCode}
\begin{english}
\begin{minted}{csharp}
public static bool IsAccepted(object[] arr)
\end{minted}
\end{english}
\end{boxCode}
    تتلقى العملية مصفوفة كائنات عامة تمثل موظفي الشركة، وتُعيد \textenglish{true} إذا قُبل الاقتراح، وإلا تُعيد \textenglish{false}. اكتب تطبيقاً كاملاً لهذه العملية.

\end{enumerate}


% ==============================================================
% الحلول النموذجية المقسمة وفق معايير المشروع لتفادي تجاوز الصفحة
% ==============================================================
\ifwithsols
\clearpage

\begin{boxSolution}[الفرع أ: مخطط الهرمية بين الفئات]
الفئة العليا الأساسية هي \textenglish{Employee}، وترث منها الفئات الثلاث الأخرى:

\begin{center}
\begin{tikzpicture}[
    classbox/.style={draw, rectangle, rounded corners=2pt, minimum width=3.4cm, minimum height=0.9cm, align=center, fill=blue!10, font=\bfseries},
    arrow/.style={-{Triangle[open]}, thick}
]
    \node[classbox] (emp) at (0, 2) {\textenglish{Employee}};
    \node[classbox] (union) at (-4, 0) {\textenglish{UnionMember}};
    \node[classbox] (tech) at (0, 0) {\textenglish{Technician}};
    \node[classbox] (sup) at (4, 0) {\textenglish{Supervisor}};

    \draw[arrow] (union.north) -- (emp.south west);
    \draw[arrow] (tech.north) -- (emp.south);
    \draw[arrow] (sup.north) -- (emp.south east);
\end{tikzpicture}
\end{center}
\end{boxSolution}

\clearpage

\begin{boxSolution}[الفئة الأساسية: \textenglish{Employee}]
{\small
\begin{english}
\begin{minted}{csharp}
public class Employee
{
    protected string id;
    protected string name;
    protected int seniority;

    public Employee(string id, string name, int seniority)
    {
        this.id = id;
        this.name = name;
        this.seniority = seniority;
    }

    public string GetId() { return this.id; }
    public string GetName() { return this.name; }
    public int GetSeniority() { return this.seniority; }

    public virtual int GetScore()
    {
        return 4 + this.seniority;
    }

    public virtual double GetSalary()
    {
        return 5000 + (this.seniority * 500);
    }

    public virtual bool YearlyBonus()
    {
        return this.seniority >= 5;
    }

    public override string ToString()
    {
        return $"[Employee] ID: {id}, Name: {name}, Seniority: {seniority}";
    }
}
\end{minted}
\end{english}
}
\end{boxSolution}

\clearpage

\begin{boxSolution}[الفئة الفرعية: \textenglish{Technician}]
{\small
\begin{english}
\begin{minted}{csharp}
public class Technician : Employee
{
    private int computers;

    public Technician(string id, string name, int seniority, int computers)
        : base(id, name, seniority)
    {
        this.computers = computers;
    }

    public int GetComputers() { return this.computers; }

    public override int GetScore()
    {
        return base.GetScore() + this.computers;
    }

    public override double GetSalary()
    {
        return base.GetSalary() + ((this.computers / 10) * 250);
    }

    public override bool YearlyBonus()
    {
        return base.YearlyBonus() && (this.computers >= 100);
    }

    public override string ToString()
    {
        return base.ToString() + $", Computers: {computers}";
    }
}
\end{minted}
\end{english}
}
\end{boxSolution}

\clearpage

\begin{boxSolution}[الفئة الفرعية: \textenglish{UnionMember}]
{\small
\begin{english}
\begin{minted}{csharp}
public class UnionMember : Employee
{
    private int unionYears;

    public UnionMember(string id, string name, int seniority, int unionYears)
        : base(id, name, seniority)
    {
        this.unionYears = unionYears;
    }

    public int GetUnionYears() { return this.unionYears; }

    public override int GetScore()
    {
        return (2 * base.GetScore()) + (2 * this.unionYears);
    }

    public override bool YearlyBonus()
    {
        return false;
    }

    public override string ToString()
    {
        return base.ToString() + $", Union Years: {unionYears}";
    }
}
\end{minted}
\end{english}
}
\end{boxSolution}

\clearpage

\begin{boxSolution}[الفئة الفرعية: \textenglish{Supervisor} (الجزء 1: التعريف والعملية البنائية والنقاط)]
{\small
\begin{english}
\begin{minted}{csharp}
public class Supervisor : Employee
{
    private Employee[] team; // Array of 15 employees

    public Supervisor(string id, string name, int seniority, Employee[] team)
        : base(id, name, seniority)
    {
        this.team = team;
    }

    public Employee[] GetTeam() { return this.team; }

    public override int GetScore()
    {
        int totalScore = base.GetScore();
        if (this.team != null)
        {
            for (int i = 0; i < this.team.Length; i++)
            {
                if (this.team[i] != null)
                {
                    totalScore += this.team[i].GetScore();
                }
            }
        }
        return totalScore;
    }
\end{minted}
\end{english}
}
\end{boxSolution}

\clearpage

\begin{boxSolution}[الفئة الفرعية: \textenglish{Supervisor} (الجزء 2: الراتب والمكافأة والطباعة)]
{\small
\begin{english}
\begin{minted}{csharp}
    public override double GetSalary()
    {
        double totalSalary = base.GetSalary();
        if (this.team != null)
        {
            for (int i = 0; i < this.team.Length; i++)
            {
                if (this.team[i] is Technician)
                {
                    totalSalary += 100;
                }
                else if (this.team[i] is UnionMember)
                {
                    totalSalary += 50;
                }
            }
        }
        return totalSalary;
    }

    public override bool YearlyBonus()
    {
        return false;
    }

    public override string ToString()
    {
        int count = (team != null) ? team.Length : 0;
        return base.ToString() + $", Supervising {count} members";
    }
}
\end{minted}
\end{english}
}
\end{boxSolution}

\clearpage

\begin{boxSolution}[الفروع هـ + و: فئة الشركة \textenglish{Company}]
العلاقة بين \textenglish{Company} و \textenglish{Employee} هي علاقة احتواء وتركيب وليست وراثة.
{\small
\begin{english}
\begin{minted}{csharp}
public class Company
{
    public const int MAX = 100;
    private Employee[] employees;
    private int current;

    public Company()
    {
        this.employees = new Employee[MAX];
        this.current = 0;
    }

    public bool AddEmployee(Employee e)
    {
        if (this.current < MAX && e != null)
        {
            this.employees[this.current] = e;
            this.current++;
            return true;
        }
        return false;
    }

    public double TotalSalary()
    {
        double total = 0;
        for (int i = 0; i < this.current; i++)
        {
            total += this.employees[i].GetSalary();
        }
        return total;
    }
}
\end{minted}
\end{english}
}
\end{boxSolution}

\clearpage

\begin{boxSolution}[الفرع ح: عملية التصويت العامة \textenglish{IsAccepted}]
المؤيدون هم الموظفون العاديون والتقنيون، بينما المعارضون هم أعضاء لجنة العمال ومسؤولو المشاريع:
{\small
\begin{english}
\begin{minted}{csharp}
public static bool IsAccepted(object[] arr)
{
    if (arr == null)
    {
        return false;
    }

    int forPoints = 0;
    int againstPoints = 0;

    for (int i = 0; i < arr.Length; i++)
    {
        if (arr[i] is Employee emp)
        {
            if (emp is UnionMember || emp is Supervisor)
            {
                againstPoints += emp.GetScore();
            }
            else // Regular Employee or Technician
            {
                forPoints += emp.GetScore();
            }
        }
    }

    return forPoints > againstPoints;
}
\end{minted}
\end{english}
}
\end{boxSolution}

\fi

\vspace{1cm}
\begin{flushleft}
أرجو لكم وقتًا ممتعًا.

الأستاذ محمود اغبارية.
\end{flushleft}

\end{document}
